\subsection*{一步一步建立直观}

\paragraph{1. 把立方体的变化拆成薄层}

切换到立方函数，把边长 $a$ 增加 $h$。三个面旁各添一层厚度 $h$ 的薄片，主项合计为 $3a^2h$；棱边附近的重叠修正含 $h^2$，角上的小块含 $h^3$。体积图像解释了为什么系数是三，而不是一个需要记住的巧合。



\paragraph{2. 先除以步长，再比较阶数}

对立方函数，差商是 $3a^2+3ah+h^2$。右图可先切回平方函数，比较其中只有一个一阶误差项的情形。两者都不是把高阶项直接删掉，而是在除以 $h$ 之后，证明剩余误差随 $h$ 趋零。然后再推广到整数次幂。



\paragraph{3. 把局部近似与大步长分开}

将步长调大，割线明显偏离切线。幂法则给出的是观察点附近的一阶倍率，不能用它精确预测任意远处的变化。对于一般实数指数，还需要说明函数定义域；本窗口的平方、立方只是几何推导最直接的两个例子。


